\section{AI 辅助使用与复现说明}
本项目使用 AI 辅助任务理解、公开实现调研、脚本示例、实验流程设计、错误排查、结果解读、图表生成与论文初稿撰写。模型训练及官方测试由学生在本机 Visual Studio 中实际运行，本文数值来自保存的运行记录。文中数值经实验日志及预测文件核对；本文文字与代码仍须由学生逐项核实并理解，提交与答辩责任由学生承担。

关键提示包括：“选择一个数据集，难度大一点的，来指导我完成大作业完成监督学习功能以及撰写专业规范的学术论文”；“你可以先去 GitHub 上找一找有没有用这个数据集来做监督学习的”；“写论文的时候，用开源的 latex 论文模板代码，注意要完成的是一篇规范的论文”。完整阶段记录见仓库中的 \texttt{har\_project/AI使用记录.md}。

代码与留档结果位于 \url{https://github.com/Gyfff000/PR_ML}。脚本 01--08 对应传统模型与 PCA，11--18 对应神经网络数据准备、试训练、分组验证与复验，19、20 分别对应最终训练和冻结模型测试。最终配置和模型保存在 \texttt{models/final}，指标及逐窗口预测保存在 \texttt{outputs/final\_test}，测试记录包含所用模型文件的 SHA-256。论文表格与图像由保存结果生成，并保留来源文件哈希清单。

排版使用开源 ElegantPaper v0.12，保留原始类文件、完整上游源码、许可证及固定提交信息。正文采用中文学术论文结构，参考文献采用 GB/T 7714 数字顺序编码样式。模板仅用于排版，不构成对本文方法或实验的学术背书。
